% VI/02 exercise solutions.

\subsection*{Solution VI.02.1}
The whole space and empty set are \(V((0))\) and \(V((1))\).  Arbitrary intersections satisfy
\[
\bigcap_\lambda V(I_\lambda)=V\!\left(\sum_\lambda I_\lambda\right),
\]
and finite unions satisfy \(V(I)\cup V(J)=V(IJ)\).  Hence the algebraic sets satisfy the
closed-set axioms.

\subsection*{Solution VI.02.2}
If \(Y\) is closed in \(X\), then \(Y=X\cap Z\) for an ambient closed set \(Z=V(J)\).
Therefore \(Y=V_X(J)\).  Conversely every \(V_X(J)=X\cap V(J)\) is closed in \(X\).

\subsection*{Solution VI.02.3}
Every \(f\in I(Y)\) vanishes on \(Y\), so \(Y\subseteq V(I(Y))\).  If \(C=V(J)\) is closed
and contains \(Y\), then \(J\subseteq I(Y)\), hence
\[
V(I(Y))\subseteq V(J)=C.
\]
Thus \(V(I(Y))\) is the smallest closed set containing \(Y\).

\subsection*{Solution VI.02.4}
For \(x\in X\),
\[
x\in D_X(f)\cap D_X(g)
\iff f(x)\neq0\text{ and }g(x)\neq0
\iff f(x)g(x)\neq0.
\]
Thus the intersection is \(D_X(fg)\).  Iterating gives
\[
\bigcap_{i=1}^rD_X(f_i)=D_X(f_1\cdots f_r).
\]

\subsection*{Solution VI.02.5}
For a field element \(a\), \(a^m=0\) iff \(a=0\).  Hence \(f(x)^m\neq0\) iff \(f(x)\neq0\),
so \(D_X(f)=D_X(f^m)\).  Taking complements also gives \(V_X(f)=V_X(f^m)\).

\subsection*{Solution VI.02.6}
A point \(x=(x_1,\dots,x_n)\) lies in the displayed zero locus exactly when
\(x_i=a_i\) for every \(i\), so the zero locus is \(\set a\).  Finite unions of closed
singletons are closed.

\subsection*{Solution VI.02.7}
If \(V(I)\) is proper, choose \(0\neq f\in I\).  Then \(V(I)\subseteq V(f)\), and \(V(f)\)
is finite.  Conversely
\[
\set{a_1,\dots,a_r}=V\!\left(\prod_i(x-a_i)\right).
\]

\subsection*{Solution VI.02.8}
Nonempty opens are cofinite.  If \(U,V\neq\varnothing\), then the complement of \(U\cap V\)
is the finite union of the finite complements of \(U\) and \(V\).  Since \(k\) is infinite,
\(U\cap V\neq\varnothing\).  Therefore disjoint nonempty neighborhoods cannot exist.

\subsection*{Solution VI.02.9}
If an infinite \(Y\) were not dense, its closure would be a proper closed subset of
\(\mathbb A_k^1\), hence finite, contradicting \(Y\subseteq\overline Y\).

\subsection*{Solution VI.02.10}
Every singleton is closed.  Since \(k^n\) is finite, every subset is a finite union of
singletons and hence closed.  Its complement is also closed, so every subset is open.

\subsection*{Solution VI.02.11}
On \(X\), \(x=0\) gives \((0,0)\), while \(x=1\) gives \((1,1)\).  Thus
\[
V_X(x)=\set{(0,0)},\qquad D_X(x)=X\setminus\set{(0,0)},
\]
\[
V_X(x-1)=\set{(1,1)},\qquad D_X(x-1)=X\setminus\set{(1,1)}.
\]

\subsection*{Solution VI.02.12}
Each polynomial \(f:k^n\to k\) is Euclidean continuous, so \(V(f)=f^{-1}(0)\) is Euclidean
closed.  Arbitrary algebraic sets are intersections of such zero sets and are therefore
Euclidean closed.  The subset \(\mathbb Z\subseteq\mathbb R\) or \(\mathbb C\) is Euclidean
closed but, being infinite, is Zariski dense in the affine line and hence is not proper
Zariski closed.

\subsection*{Solution VI.02.13}
By the closure formula, \(Y\) is dense iff \(V(I(Y))=\mathbb A_k^n\).  Over an infinite
field, the only polynomial vanishing on every point of affine space is zero.  Therefore
density is equivalent to \(I(Y)=(0)\).

\subsection*{Solution VI.02.14}
Induct on \(n\).  Let \(f\) vanish on the product and write
\[
f=\sum_{j=0}^d a_j(x_1,\dots,x_{n-1})x_n^j.
\]
For each point of \(A_1\times\cdots\times A_{n-1}\), the resulting one-variable polynomial
vanishes on the infinite set \(A_n\), so every coefficient value is zero.  By induction each
\(a_j\) is zero.  Thus \(f=0\), proving density.

\subsection*{Solution VI.02.15}
By De Morgan's law,
\[
D_X(f)\cup D_X(g)=X\setminus\bigl(V_X(f)\cap V_X(g)\bigr)
=X\setminus V_X(f,g).
\]
A complement of a two-equation zero set is open, but the formula does not express it as
\(D_X(h)\) for one polynomial \(h\).

\subsection*{Solution VI.02.16}
Write
\[
\mathbb A^1\setminus U=\set{a_1,\dots,a_r}=V(f(x)),
\qquad
\mathbb A^1\setminus V=\set{b_1,\dots,b_s}=V(g(y)).
\]
Then
\[
U\times V=\set{(x,y):f(x)\neq0,\ g(y)\neq0}=D(f(x)g(y)),
\]
which is Zariski open in \(\mathbb A_k^2\).

\subsection*{Solution VI.02.17}
A matrix has rank at most \(r\) exactly when every \((r+1)\times(r+1)\) minor vanishes.
Each minor is a polynomial in the matrix entries, so the rank-at-most-\(r\) locus is their
common zero locus and hence is Zariski closed.

\subsection*{Solution VI.02.18}
A point lies in \((k^\times)^n\) exactly when all coordinates are nonzero, equivalently when
\(x_1\cdots x_n\neq0\).  Hence the set is \(D(x_1\cdots x_n)\).  The defining polynomial is
nonzero; over an infinite field the proposition on principal opens makes this set dense.
